%HOMEWORK template for M 325K
\documentclass{article}
\headheight=7pt         \topmargin=14pt
\textheight=574pt       \textwidth=445pt
\oddsidemargin=18pt     \evensidemargin=18pt

%Begin package import

\usepackage{amsmath, amssymb, amsthm}
\usepackage{mathtools}
\usepackage{pdfpages}
\usepackage{tikz-cd}

%End package import
%Begin custom commands

\newcommand{\C}{\mathbb{C}}
\newcommand{\R}{\mathbb{R}}
\newcommand{\Q}{\mathbb{Q}}
\newcommand{\Z}{\mathbb{Z}}
\newcommand{\N}{\mathbb{N}}
\newcommand{\abs}[1]{\lvert #1\rvert}
\newcommand{\RM}[1]{\mathrm{#1}}
\newcommand{\BF}[1]{\textbf{#1}}
\newcommand{\e}{\epsilon}
\newcommand{\ceil}[1]{\lceil{#1}\rceil}
\newcommand{\floor}[1]{\lfloor{#1}\rfloor}
\newcommand{\rarr}[0]{\rightarrow}
\newcommand{\IM}[1]{\text{Im}(#1)}
\newcommand{\Hom}[0]{\text{Hom}}
\newcommand{\Pow}[1]{\text{Pow}(#1)}

%End custom commands
%Begin custom theorems

\newtheorem{innercustomgeneric}{\customgenericname}
\providecommand{\customgenericname}{}
\newcommand{\newcustomtheorem}[2]{%
\newenvironment{#1}[1]
{%
\renewcommand\customgenericname{#2}%
\renewcommand\theinnercustomgeneric{##1}%
\innercustomgeneric%
}
{\endinnercustomgeneric}
}

\newcustomtheorem{THRM}{Theorem}
\newcustomtheorem{LEM}{Lemma}
\newcustomtheorem{EX}{Exercise}
\newcustomtheorem{COR}{Corollary}
\newcustomtheorem{PROB}{Problem}
\newcustomtheorem{claim}{Claim}

\newenvironment{solution}
{\renewcommand\qedsymbol{$\blacksquare$}\begin{proof}[Solution]}
{\end{proof}}

\newenvironment{definition}
{\renewcommand\qedsymbol{$\blacksquare$}\begin{proof}[Definition]}
{\end{proof}}

%End custom theorems
\begin{document}

\title{Homework 1}
\author{Arham Lodha}%put your name here
\date{January 16, 2024}%enter the due date here
\maketitle

\begin{PROB}{1}\label{Problem 1.}
    Prove the Distributive Laws:
    
    \begin{enumerate}
        \item $A \cap (B \cup C) = (A \cap B) \cup (A \cap C)$
        \item $A \cup (B \cap C) = (A \cup B) \cap (A \cup C)$ 
    \end{enumerate}
\end{PROB}
\begin{proof}
    \begin{enumerate}
        \item $\forall x \in A \cap (B \cup C)$. By definition of a intersection, $x \in A$ and $x \in B \cup C$. By definition of union, $x \in B$ or $x \in C$. If $x \in B$, by definition of intersection, $x \in A \cap B$. If $x \in C$, by definition of intersection, $x \in A \cap C$. By definition of union, $x \in (A \cap B) \cup (A \cap C)$. Thus $A \cap (B \cup C) \subseteq (A \cap B) \cup (A \cap C)$. $\forall x \in (A \cap B) \cup (A \cap C)$, by definition of union, $x \in A \cap B$ or $x \in A \cap C$. By definition of intersection, $x \in A$ and $x \in B$ or $x \in C$. By definition of union, $x \in B \cup C$. Thus $x \in A$ and $x \in B \cup C$, by definition of intersection, $x \in A \cap (B \cup C)$. Thus $(A \cap B) \cup (A \cap C) \subseteq A \cap (B \cup C)$. Thus $(A \cap B) \cup (A \cap C) = A \cap (B \cup C)$.
        

        \item $\forall x \in A \cup (B \cap C)$. By definition of union, $x \in A$ or $x \in B \cap C$. If $x \in A$, by definition of union, $x \in A \cup B$ and $x \in A \cup C$. If $x \in B \cap C$, then $x \in B$ and $x \in C$. By definition of union, $x \in A \cup B$ and $x \in A \cup C$. By definition of intersection, $x \in (A \cup B) \cap (A \cup C)$. Thus $A \cup (B \cap C) \subseteq (A \cup B) \cap (A \cup C)$. $\forall x \in (A \cup B) \cap (A \cup C)$. By definition of intersection, $x \in A \cup B$ and $x \in A \cup C$. By definition of union, since $x \in A \cup B$, $x \in A$ or $x \in B$. Similarly, $x \in A$ or $x \in C$. If $x \in A$, by definition of union, $x \in A \cup (B \cap C)$. If $x \not \in A$, $x \in B$ and $x \in C$, then $x \in B \cap C$. Thus by definition of union, $x \in A \cup (B \cap C)$. Thus $(A \cup B) \cap (A \cup C) \subseteq A \cup (B \cap C)$. Thus $A \cup (B \cap C) = (A \cup B) \cap (A \cup C)$.          
    \end{enumerate}
    
\end{proof}

\bigskip

\begin{PROB}{2}\label{Problem 2.}
    2. Let $f(x)=x^2$ for $x \in \mathbb{R}$, and let $E=\{x \in \mathbb{R}:-1 \leq x \leq 0\}$ and $F=\{x \in \mathbb{R}: 0 \leq x \leq 1\}$.

    \begin{enumerate}
        \item Show that $E \cap F=\{0\}$, and $f(E \cap F)=\{0\}$, while $f(E)=f(F)=\{y \in \mathbb{R}: 0 \leq$ $y \leq 1\}$. Hence $f(E \cap F)$ is a proper subset of $f(E) \cap f(F)$.
        \item Find $f^{-1}(E)$ and $f^{-1}(F)$, the inverse image of $E$, and the inverse image of $F$.
    \end{enumerate}
\end{PROB}
\begin{solution}
    \begin{enumerate}
        \item \begin{proof}
            $0 \in \R$ and $-1 \leq 0 \leq 0$ and $0 \leq 0 \leq 1$. So $0 \in E$ and $0 \in F$ so by definition of intersection, $E \cap F$. So $\left\{0\right\} \subseteq E \cap F$.      $\forall x \in E \cap F$, by definition of intersection, $x \in E$ and $x \in F$. Suppose $x > 0$, then $0 < x \leq 1$. Thus $x \not \leq 0$, thus $x \not \in F$, thus $x \not \in E \cap F$. Similarly suppose, $x < 0$, then $-1 \leq x < 0$. Thus $x \not \geq 0$, thus $x \not \in F$ and $x \not \in E \cap F$. Thus $E \cap F  = \emptyset$ or $E \cap F = \{0\}$. Thus $E \cap F \subseteq \left\{0\right\}$. Thus $E \cap F = \left\{0\right\}$. 

            By definition of image, $f(E \cap F) = f(\left\{0\right\}) = \left\{f(x) : x \in E \cap F\right\} = \left\{f(0)\right\} = \left\{0\right\}$. $f(E) = \left\{x^2 : -1 \leq x \leq 0 \right\}$ and $f(F) = \left\{ x^2 : 0 \leq x \leq 1 \right\}$. Both sets represent non-negative numbers between 0 and 1. So $f(E) = f(F) = \left\{ y : 0 \leq y \leq 1 \right\}$. $f(F) = f(E) \implies f(E) \cap f(F) = f(E) = f(F)$. Since $0 \in f(F) \implies 0 \in f(E) \cap f(F) \implies f(E \cap F) \subset f(E) \cap f(F)$. 
        \end{proof}

        \item $\forall x \in E \setminus \left\{ 0 \right\}, \not \exists a \in \R \ni f(a) = a^2 = x $, since $f$ is strictly greater than or equal to 0 for all real numbers. For $x = 0$, only 0 is such a value such that $0^2 = 0$, because the square of any nonzero number is strictly positive. So $f^{-1}(E) = \left\{ 0 \right\}$.
        
        $\forall x \in F$, by definition $0 \leq x \leq 1$. For any positive real number $x, \exists \sqrt{x} \ni \sqrt{x}^2 = x$ and $(-\sqrt{x})^2 = x$. The square root of a number between 0 and 1 is also between 0 and 1. So $\exists \pm\sqrt{x} \in [-1, 1]$. So $f^{-1}(F) \subseteq [-1, 1]$. $\forall y \in [-1, 1]$, $f(y) = y^2 \in F$, because the square of a real number between -1 and 1 is also between 0 and 1. So $y \in f^{-1}(F)$. So $[-1, 1] \subseteq f^{-1}(F) \implies f^{-1}(F) = [-1, 1]$.           
    \end{enumerate}
    
\end{solution}


\bigskip

\begin{PROB}{3}\label{Problem 3.}
    Show that if $f: A \rightarrow B$, and $G, H$ are subsets of $B$, then:

    \begin{enumerate}
        \item $f^{-1}(G \cup H)=f^{-1}(G) \cup f^{-1}(H)$.
        \item $f^{-1}(G \cap H)=f^{-1}(G) \cap f^{-1}(H)$.
    \end{enumerate}
\end{PROB}
\begin{solution}
    \begin{enumerate}
        \item $\forall x \in f^{-1}(G \cup H)$, by definition of inverse image, $f(x) \in G \cup H \implies f(x) \in G$ or $f(x) \in H$. If $f(x) \in G \implies x \in f^{-1}(G)$, by definition of union, $x \in f^{-1}(G) \cup f^{-1}(H)$. Similarly, if $f(x) \in \implies f^{-1}(H)$, by definition of union, $x \in f^{-1}(G) \cup f^{-1}(H)$. So $f^{-1}(G \cup H) \subseteq f^{-1}(G) \cup f^{-1}(H)$. 
        
        $\forall x \in f^{-1}(G) \cup f^{-1}(H)$, by definition of union, $x \in f^{-1}(G)$ or $x \in f^{-1}(H)$. Thus $f(x) \in G$ or $f(x) \in H$ and by definition of union, $f(x) \in G \cup H$. By definition of inverse image. $x \in f^{-1}(G \cup H)$. Thus $f^{-1}(G) \cup f^{-1}(H) \subseteq f^{-1}(G \cup H) \implies f^{-1}(G \cup H) = f^{-1}(G) \cup f^{-1}(H)$
        
        \item $\forall x \in f^{-1}(G \cap H)$, by definition of the inverse image, $f(x) \in G \cap H$. By definition of intersection, $f(x) \in G$ and $f(x) \in H$. By definition of inverse image, $x \in f^{-1}(G)$ and $x \in f^{-1}(H)$. By definition of intersection, $x \in f^{-1}(G) \cap f^{-1}(H)$. Thus $f^{-1}(G \cap H) \subseteq f^{-1}(G) \cap f^{-1}(H)$.
        
        $\forall x \in f^{-1}(G) \cap f^{-1}(H)$, by definition of intersection, $x \in f^{-1}(G)$ and $x \in f^{-1}(H)$. By definition of inverse image, $f(x) \in G$ and $f(x) \in H$. By definition of intersection, $f(x) \in G \cap H$. By definition of inverse image, $x \in f^{-1}(G \cap H)$. Thus $f^{-1}(G) \cap f^{-1}(H) \subseteq f^{-1}(G \cap H)$. Thus $f^{-1}(G \cap H)=f^{-1}(G) \cap f^{-1}(H)$.          
    \end{enumerate}
\end{solution}

\bigskip

\begin{PROB}{4}\label{Problem 4.}
    Prove that if $f : A \to B$ is bijective and $g : B \to C$ is bijective, then the composite $g \circ f$ is a bijective map of A onto C.
\end{PROB}
\begin{proof}

    \textbf{Injectivity:} $\forall x, y \in A \ni (g \circ f)(x) = (g \circ f)(x) \implies g(f(x)) = g(f(y))$. By definition of bijective functions, g is injective. By definition of injective, since $g(f(x)) = g(f(y)) \implies f(x) = f(y)$. By definition of bijective functions, f is injective. By definition of injective functions, since $f(x) = f(y) \implies x = y$. Thus $g \circ f$ is injective.
    
    \textbf{Surjectivity}: By definition of bijection, g and f are surjective. $\forall y \in C$, since g is surjective, $\exists b \in B \ni g(b) = y$. Since $f$ is surjective, $\exists x \in A \ni f(x) = b$. Thus $g(f(x)) = g(b) = y$. Thus $(g \circ f)(x) = y$. Thus $g \circ f  $ is a surjective function. 
    
    Thus by definition of bijective function, $g \circ f$ is a bijective map. 
\end{proof}

\end{document}